%% tikzlibrarycommtikz.base.code.tex
%% Core keys, PGF drawing helpers, multi-port anchor engine, and generic block shapes.

% ==============================================================================
% 1. Key Definitions (/tikz/commtikz and /tikz forwarding keys)
% ==============================================================================
\pgfkeys{
  /tikz/commtikz/.is family,
  /tikz/commtikz/.cd,
  unit/.initial=1cm,
  block scale/.initial=0.85,
  circuit scale/.initial=1.10,
  component scale/.initial=1,
  arrow length/.initial=0.10,
  arrow width/.initial=0.12,
  arrow nominal width/.initial=1,
  arrow nominal height/.initial=1,
  fixed W/.initial={},
  fixed E/.initial={},
  width/.initial=1,
  height/.initial=1,
  height1/.initial=2.2,
  height2/.initial=1.0,
  radius/.initial=0.5,
  bend radius/.initial=1.0,
  numN/.initial=1,
  numS/.initial=1,
  numE/.initial=1,
  numW/.initial=1,
  port span x/.initial={},
  port span y/.initial={},
  port order/.initial=1, % 1 = top-to-bottom (default), -1 = bottom-to-top
  line width/.initial=0.6pt,
  inner line width/.initial=0.45pt,
  signal line width/.initial=0.6pt,
  label distance/.initial=3pt,
  fill/.initial=white,
  comp label/.initial={},
  sublabel I/.initial={I},
  sublabel Q/.initial={Q},
  response/.initial={},
  tuneup/.initial=1,
  flip/.initial=0,
  ccw/.initial=0,
  electrodes/.initial=0,
  add drop/.initial=0,
  show phase/.initial=0,
  boxed/.initial=0,
  tunable/.initial=0,
  sector N/.initial={},
  sector S/.initial={},
  sector E/.initial={},
  sector W/.initial={},
  % Port count group helper: ports={W=2, E=4, N=1, S=1}
  ports/.is family,
  ports/N/.style={/tikz/commtikz/numN=#1},
  ports/S/.style={/tikz/commtikz/numS=#1},
  ports/E/.style={/tikz/commtikz/numE=#1},
  ports/W/.style={/tikz/commtikz/numW=#1},
}

% Top-level /tikz convenience keys so users can pass them directly on nodes
\tikzset{
  commtikz block scale/.style={/tikz/commtikz/block scale=#1},
  commtikz circuit scale/.style={/tikz/commtikz/circuit scale=#1},
  commtikz legacy sizes/.style={commtikz block scale=1, commtikz circuit scale=1},
  commtikz circuit node/.style={
    /tikz/commtikz/component scale=\pgfkeysvalueof{/tikz/commtikz/circuit scale},
    /tikz/commtikz/arrow length=0.065,
    /tikz/commtikz/arrow width=0.080,
  },
  commtikz arrow reference/.style args={#1/#2}{
    /tikz/commtikz/arrow nominal width=#1,
    /tikz/commtikz/arrow nominal height=#2,
  },
  commtikz fixed ports/.style args={#1/#2}{
    /tikz/commtikz/fixed W=#1, /tikz/commtikz/fixed E=#2,
  },
  commtikz unit/.style={/tikz/commtikz/unit=#1},
  commtikz scale/.style={/tikz/commtikz/unit=\the\dimexpr#1\dimexpr1cm\relax\relax},
  commtikz line width/.style={
    /tikz/commtikz/line width=#1,
    /tikz/commtikz/signal line width=#1,
    /tikz/commtikz/inner line width/.expanded={\the\dimexpr0.75\dimexpr#1\relax\relax}
  },
  commtikz block line width/.style={/tikz/commtikz/line width=#1},
  commtikz inner line width/.style={/tikz/commtikz/inner line width=#1},
  commtikz signal line width/.style={/tikz/commtikz/signal line width=#1},
  commtikz label distance/.style={/tikz/commtikz/label distance=#1},
  commtikz match font/.style={commtikz line width=0.055em},
  width/.style={/tikz/commtikz/width=#1},
  height/.style={/tikz/commtikz/height=#1},
  height1/.style={/tikz/commtikz/height1=#1},
  height2/.style={/tikz/commtikz/height2=#1},
  block width/.style={/tikz/commtikz/width=#1},
  block height/.style={/tikz/commtikz/height=#1},
  comp radius/.style={/tikz/commtikz/radius=#1},
  bend radius/.style={/tikz/commtikz/bend radius=#1},
  numN/.style={/tikz/commtikz/numN=#1},
  numS/.style={/tikz/commtikz/numS=#1},
  numE/.style={/tikz/commtikz/numE=#1},
  numW/.style={/tikz/commtikz/numW=#1},
  port span x/.style={/tikz/commtikz/port span x={#1}},
  port span y/.style={/tikz/commtikz/port span y={#1}},
  port span/.style={/tikz/commtikz/port span x={#1}, /tikz/commtikz/port span y={#1}},
  ports/.style={/tikz/commtikz/ports/.cd,#1,/tikz/.cd},
  port order/.is choice,
  port order/top to bottom/.style={/tikz/commtikz/port order=1},
  port order/bottom to top/.style={/tikz/commtikz/port order=-1},
  fillcol/.style={fill=#1,/tikz/commtikz/fill=#1},
  comp label/.style={/tikz/commtikz/comp label={#1}},
  sublabel I/.style={/tikz/commtikz/sublabel I={#1}},
  sublabel Q/.style={/tikz/commtikz/sublabel Q={#1}},
  response/.style={/tikz/commtikz/response={#1}},
  tuneup/.is choice,
  tuneup/true/.style={/tikz/commtikz/tuneup=1},
  tuneup/false/.style={/tikz/commtikz/tuneup=0},
  tuneup/.default=true,
  mirror flip/.is choice,
  mirror flip/true/.style={/tikz/commtikz/flip=1},
  mirror flip/false/.style={/tikz/commtikz/flip=0},
  mirror flip/.default=true,
  mux flip/.is choice,
  mux flip/true/.style={/tikz/commtikz/flip=1},
  mux flip/false/.style={/tikz/commtikz/flip=0},
  mux flip/.default=true,
  bs flip/.is choice,
  bs flip/true/.style={/tikz/commtikz/flip=1},
  bs flip/false/.style={/tikz/commtikz/flip=0},
  bs flip/.default=true,
  ccw/.is choice,
  ccw/true/.style={/tikz/commtikz/ccw=1},
  ccw/false/.style={/tikz/commtikz/ccw=0},
  ccw/.default=true,
  electrodes/.is choice,
  electrodes/true/.style={/tikz/commtikz/electrodes=1},
  electrodes/false/.style={/tikz/commtikz/electrodes=0},
  electrodes/.default=true,
  add drop/.is choice,
  add drop/true/.style={/tikz/commtikz/add drop=1},
  add drop/false/.style={/tikz/commtikz/add drop=0},
  add drop/.default=true,
  show phase/.is choice,
  show phase/true/.style={/tikz/commtikz/show phase=1},
  show phase/false/.style={/tikz/commtikz/show phase=0},
  show phase/.default=true,
  boxed/.is choice,
  boxed/true/.style={/tikz/commtikz/boxed=1},
  boxed/false/.style={/tikz/commtikz/boxed=0},
  boxed/.default=true,
  tunable/.is choice,
  tunable/true/.style={/tikz/commtikz/tunable=1},
  tunable/false/.style={/tikz/commtikz/tunable=0},
  tunable/.default=true,
  sector N/.style={/tikz/commtikz/sector N={#1}},
  sector S/.style={/tikz/commtikz/sector S={#1}},
  sector E/.style={/tikz/commtikz/sector E={#1}},
  sector W/.style={/tikz/commtikz/sector W={#1}},
  % Base visual style applied to all commtikz block nodes
  commtikz base node/.style={
    /tikz/commtikz/component scale=\pgfkeysvalueof{/tikz/commtikz/block scale},
    draw,
    fill=\pgfkeysvalueof{/tikz/commtikz/fill},
    line width=\pgfkeysvalueof{/tikz/commtikz/line width},
    line cap=butt,
    line join=miter,
    inner sep=0pt,
    outer sep=0pt,
  },
  commtikz bare node/.style={
    /tikz/commtikz/component scale=\pgfkeysvalueof{/tikz/commtikz/block scale},
    draw,
    line width=\pgfkeysvalueof{/tikz/commtikz/line width},
    line cap=butt,
    line join=miter,
    inner sep=0pt,
    outer sep=0pt,
  },
  % Shared default footprint for ordinary square building blocks.
  commtikz square node/.style={commtikz base node, width=1, height=1},
}

% ==============================================================================
% 2. Low-level PGF Helper Macros and Dimension Registers for Shapes
% ==============================================================================
% Labels are created after the component's option scope has ended. Save the
% clearance on the parent shape and apply it only to that shape's labels.
\def\commtikz@savelabeldistance{%
  \edef\commtikz@saved@labeldistance{\pgfkeysvalueof{/tikz/commtikz/label distance}}%
  \addtosavedmacro\commtikz@saved@labeldistance
}
\def\commtikz@labeldefaults{%
  \pgfutil@ifundefined{pgf@sh@ma@\tikzlastnode}{}{%
    \begingroup
      \let\commtikz@savedlabel\pgfutil@empty
      \let\commtikz@saved@labeldistance\pgfutil@empty
      \csname pgf@sh@ma@\tikzlastnode\endcsname
      \commtikz@savedlabel
      \ifx\commtikz@saved@labeldistance\pgfutil@empty
        \endgroup
      \else
        \edef\commtikz@next{\endgroup
          \noexpand\tikzset{label distance=\commtikz@saved@labeldistance}}%
        \commtikz@next
      \fi
  }%
}
\tikzset{every label/.append code={\commtikz@labeldefaults}}

\newdimen\commtikz@u
\newdimen\commtikz@hw
\newdimen\commtikz@hh
\newdimen\commtikz@phw
\newdimen\commtikz@phh
\newdimen\commtikz@hha
\newdimen\commtikz@hhb
\newdimen\commtikz@hhW
\newdimen\commtikz@hhE
\newdimen\commtikz@r
\newdimen\commtikz@rk
\newdimen\commtikz@br
\newdimen\commtikz@bdx
\newdimen\commtikz@bdy
\newdimen\commtikz@ndx
\newdimen\commtikz@ndy
\newdimen\commtikz@fr
\newdimen\commtikz@fy
\newdimen\commtikz@fdx
\newdimen\commtikz@pcr
\newdimen\commtikz@pcy
\newdimen\commtikz@pcdx
\newdimen\commtikz@voar
\newdimen\commtikz@xa
\newdimen\commtikz@xb
\newdimen\commtikz@xc
\newdimen\commtikz@xd
\newdimen\commtikz@xe
\newdimen\commtikz@xf
\newdimen\commtikz@ya
\newdimen\commtikz@xO
\newdimen\commtikz@xA
\newdimen\commtikz@xB
\newdimen\commtikz@xC
\newdimen\commtikz@xD
\newdimen\commtikz@yM
\newdimen\commtikz@yH
\newdimen\commtikz@yL
\newdimen\commtikz@fbggh
\newdimen\commtikz@gx
\newdimen\commtikz@ar
\newdimen\commtikz@rr
\newdimen\commtikz@rx
\newdimen\commtikz@ry
\newdimen\commtikz@sy
\newdimen\commtikz@arrowunit
\newdimen\commtikz@arr@xa
\newdimen\commtikz@arr@ya
\newdimen\commtikz@arr@xb
\newdimen\commtikz@arr@yb
\newdimen\commtikz@arr@hl
\newdimen\commtikz@arr@hw

% Use current stroke color as fill color inside a \pgfscope
\def\commtikz@use@strokecolor@as@fill{%
  \ifx\tikz@strokecolor\pgfutil@empty
    \pgfsetfillcolor{.}%
  \else
    \pgfsetfillcolor{\tikz@strokecolor}%
  \fi
}

% Set both PDF fill color and LaTeX text color for internal shape labels
\def\commtikz@set@labelcolor{%
  \ifx\tikz@textcolor\pgfutil@empty
    \ifx\tikz@strokecolor\pgfutil@empty
      \pgfsetfillcolor{.}\color{.}%
    \else
      \pgfsetfillcolor{\tikz@strokecolor}\color{\tikz@strokecolor}%
    \fi
  \else
    \pgfsetfillcolor{\tikz@textcolor}\color{\tikz@textcolor}%
  \fi
}

% Shared open heads use the component's normalized size, never shaft length.
% Thus e.g. NMOS and NPN heads match despite differently oriented shafts.
\def\commtikz@arrowgeometry#1#2#3#4{%
  \commtikz@arr@xa=\dimexpr#1\relax
  \commtikz@arr@ya=\dimexpr#2\relax
  \commtikz@arr@xb=\dimexpr#3\relax
  \commtikz@arr@yb=\dimexpr#4\relax
  \commtikz@arr@hl=\pgfkeysvalueof{/tikz/commtikz/arrow length}\commtikz@arrowunit\relax
  \commtikz@arr@hw=\pgfkeysvalueof{/tikz/commtikz/arrow width}\commtikz@arrowunit\relax
  \pgfpointnormalised{\pgfpoint{\dimexpr\commtikz@arr@xb-\commtikz@arr@xa\relax}{\dimexpr\commtikz@arr@yb-\commtikz@arr@ya\relax}}%
  \edef\commtikz@ux{\pgf@sys@tonumber{\pgf@x}}%
  \edef\commtikz@uy{\pgf@sys@tonumber{\pgf@y}}%
  \commtikz@bdx=\commtikz@ux\commtikz@arr@hl\relax
  \commtikz@bdy=\commtikz@uy\commtikz@arr@hl\relax
  \commtikz@ndx=-0.5\dimexpr\commtikz@uy\commtikz@arr@hw\relax
  \commtikz@ndy=0.5\dimexpr\commtikz@ux\commtikz@arr@hw\relax
}
\def\commtikz@openheadpath{%
  \pgfpathmoveto{\pgfpoint{\dimexpr\commtikz@arr@xb-\commtikz@bdx+\commtikz@ndx\relax}{\dimexpr\commtikz@arr@yb-\commtikz@bdy+\commtikz@ndy\relax}}%
  \pgfpathlineto{\pgfpoint{\commtikz@arr@xb}{\commtikz@arr@yb}}%
  \pgfpathlineto{\pgfpoint{\dimexpr\commtikz@arr@xb-\commtikz@bdx-\commtikz@ndx\relax}{\dimexpr\commtikz@arr@yb-\commtikz@bdy-\commtikz@ndy\relax}}%
}
\def\commtikz@drawopenarrow#1#2#3#4{%
  \pgfscope
    \commtikz@arrowgeometry{#1}{#2}{#3}{#4}%
    \pgfpathmoveto{\pgfpoint{\commtikz@arr@xa}{\commtikz@arr@ya}}%
    \pgfpathlineto{\pgfpoint{\commtikz@arr@xb}{\commtikz@arr@yb}}%
    \commtikz@openheadpath
    \pgfusepath{stroke}%
  \endpgfscope
}
% Head only, for an existing straight lead or tangent of a curved arrow.
\def\commtikz@drawopenhead#1#2#3#4{%
  \pgfscope
    \commtikz@arrowgeometry{#1}{#2}{#3}{#4}%
    \commtikz@openheadpath
    \pgfusepath{stroke}%
  \endpgfscope
}

\def\commtikz@setunit{%
  \pgfmathsetlength\commtikz@u{\pgfkeysvalueof{/tikz/commtikz/unit} * \pgfkeysvalueof{/tikz/commtikz/component scale}}%
}
\def\commtikz@savearrowunit{%
  \pgfmathsetlength\commtikz@arrowunit{min(2*\commtikz@hw/\pgfkeysvalueof{/tikz/commtikz/arrow nominal width},2*\commtikz@hh/\pgfkeysvalueof{/tikz/commtikz/arrow nominal height})}%
  \edef\commtikz@saved@arrowunit{\the\commtikz@arrowunit}%
  \addtosavedmacro\commtikz@saved@arrowunit
}
% Validate the final options, so later numW/ports/port-order overrides cannot
% silently move anchors away from the fixed internal waveguide topology.
\def\commtikz@validatefixedports{%
  \pgfkeysgetvalue{/tikz/commtikz/fixed W}{\commtikz@fixedW}%
  \pgfkeysgetvalue{/tikz/commtikz/fixed E}{\commtikz@fixedE}%
  \def\commtikz@saved@fixedports{0}%
  \ifx\commtikz@fixedW\pgfutil@empty\else
    \def\commtikz@saved@fixedports{1}%
    \ifnum\pgfkeysvalueof{/tikz/commtikz/numW}=\commtikz@fixedW\relax\else
      \PackageError{commtikz}{Fixed topology requires numW=\commtikz@fixedW}{Use the documented port counts for this component.}%
    \fi
    \ifnum\pgfkeysvalueof{/tikz/commtikz/numE}=\commtikz@fixedE\relax\else
      \PackageError{commtikz}{Fixed topology requires numE=\commtikz@fixedE}{Use the documented port counts for this component.}%
    \fi
    \ifnum\pgfkeysvalueof{/tikz/commtikz/port order}=1\relax\else
      \PackageError{commtikz}{Fixed topology requires top-to-bottom port order}{Use port order=top to bottom.}%
    \fi
    \pgfmathsetlength\pgf@x{abs(\commtikz@phh-\commtikz@hh)}%
    \ifdim\pgf@x>0.01pt
      \PackageError{commtikz}{Fixed topology requires full-height port span}{Remove port span y or set it equal to the block height.}%
    \fi
  \fi
  \addtosavedmacro\commtikz@saved@fixedports
}

\def\commtikz@checkportindex#1#2{%
  \ifnum\commtikz@saved@fixedports=1\relax
    \ifnum#2<\csname commtikz@num#1\endcsname\relax\else
      \PackageError{commtikz}{Port #1#2 is outside this fixed topology}{Use the documented ports for this component.}%
    \fi
  \fi
}

% Draw default component label at (#1,#2) if node text box is empty
\def\commtikz@drawcomplabel#1#2{%
  \ifdim\wd\pgfnodeparttextbox=0pt\relax
    \pgfkeysgetvalue{/tikz/commtikz/comp label}{\commtikz@lbl}%
    \ifx\commtikz@lbl\pgfutil@empty\else
      \pgfscope
        \commtikz@set@labelcolor
        \pgftext[at=\pgfpoint{#1}{#2}]{\commtikz@set@labelcolor\tikz@textfont\commtikz@lbl}%
      \endpgfscope
    \fi
  \fi
}

% Generate N0..N15, S0..S15, E0..E15, W0..W15 and 1-based north/south/east/west aliases
\def\commtikz@defrectportanchors{%
  \pgfutil@for\commtikz@idx:=0,1,2,3,4,5,6,7,8,9,10,11,12,13,14,15\do{%
    \expandafter\commtikz@defrectportanchor@i\expandafter{\commtikz@idx}%
  }%
}

\def\commtikz@defrectportanchor@i#1{%
  \anchor{N#1}{%
    \commtikz@rectparams
    \commtikz@checkportindex{N}{#1}%
    \pgfmathsetlength\pgf@x{-\commtikz@phw + (#1 + 1) * (2 * \commtikz@phw) / (\commtikz@numN + 1)}%
    \pgf@y=\commtikz@hh\relax
  }%
  \anchor{S#1}{%
    \commtikz@rectparams
    \commtikz@checkportindex{S}{#1}%
    \pgfmathsetlength\pgf@x{-\commtikz@phw + (#1 + 1) * (2 * \commtikz@phw) / (\commtikz@numS + 1)}%
    \pgf@y=-\commtikz@hh\relax
  }%
  \anchor{E#1}{%
    \commtikz@rectparams
    \commtikz@checkportindex{E}{#1}%
    \pgf@x=\commtikz@hw\relax
    \pgfmathsetlength\pgf@y{\commtikz@portorder * (\commtikz@phh - (#1 + 1) * (2 * \commtikz@phh) / (\commtikz@numE + 1))}%
  }%
  \anchor{W#1}{%
    \commtikz@rectparams
    \commtikz@checkportindex{W}{#1}%
    \pgf@x=-\commtikz@hw\relax
    \pgfmathsetlength\pgf@y{\commtikz@portorder * (\commtikz@phh - (#1 + 1) * (2 * \commtikz@phh) / (\commtikz@numW + 1))}%
  }%
}

% Inherit all rectangular anchors from #1 (usually commrect)
\def\commtikz@inheritrectanchors#1{%
  \inheritsavedanchors[from=#1]%
  \inheritanchorborder[from=#1]%
  \inheritanchor[from=#1]{center}%
  \inheritanchor[from=#1]{mid}%
  \inheritanchor[from=#1]{base}%
  \inheritanchor[from=#1]{north}%
  \inheritanchor[from=#1]{south}%
  \inheritanchor[from=#1]{east}%
  \inheritanchor[from=#1]{west}%
  \inheritanchor[from=#1]{north east}%
  \inheritanchor[from=#1]{north west}%
  \inheritanchor[from=#1]{south east}%
  \inheritanchor[from=#1]{south west}%
  \inheritanchor[from=#1]{N}%
  \inheritanchor[from=#1]{S}%
  \inheritanchor[from=#1]{E}%
  \inheritanchor[from=#1]{W}%
  \inheritanchor[from=#1]{NE}%
  \inheritanchor[from=#1]{NW}%
  \inheritanchor[from=#1]{SE}%
  \inheritanchor[from=#1]{SW}%
  \inheritanchor[from=#1]{in}%
  \inheritanchor[from=#1]{out}%
  \inheritanchor[from=#1]{input}%
  \inheritanchor[from=#1]{output}%
  \pgfutil@for\commtikz@idx:=0,1,2,3,4,5,6,7,8,9,10,11,12,13,14,15\do{%
    \expandafter\commtikz@inheritport@i\expandafter{\commtikz@idx}{#1}%
  }%
}

\def\commtikz@inheritport@i#1#2{%
  \inheritanchor[from=#2]{N#1}%
  \inheritanchor[from=#2]{S#1}%
  \inheritanchor[from=#2]{E#1}%
  \inheritanchor[from=#2]{W#1}%
}

% ==============================================================================
% 3. Base Shape: commrect
% ==============================================================================
\def\commtikz@rectparams{%
  \commtikz@savedrect
  \commtikz@arrowunit=\commtikz@saved@arrowunit\relax
  \commtikz@u=\commtikz@saved@u\relax
  \commtikz@hw=\commtikz@saved@hw\relax
  \commtikz@hh=\commtikz@saved@hh\relax
  \commtikz@phw=\commtikz@saved@phw\relax
  \commtikz@phh=\commtikz@saved@phh\relax
}

\pgfdeclareshape{commrect}{
  \savedmacro\commtikz@savedlabel{\commtikz@savelabeldistance}
  \savedmacro\commtikz@savedrect{%
    \commtikz@setunit
    \pgfmathparse{\pgfkeysvalueof{/tikz/commtikz/width}}%
    \ifpgfmathunitsdeclared
      \commtikz@hw=0.5\dimexpr\pgfmathresult pt\relax
    \else
      \commtikz@hw=0.5\dimexpr\pgfmathresult\commtikz@u\relax
    \fi
    \pgfmathparse{\pgfkeysvalueof{/tikz/commtikz/height}}%
    \ifpgfmathunitsdeclared
      \commtikz@hh=0.5\dimexpr\pgfmathresult pt\relax
    \else
      \commtikz@hh=0.5\dimexpr\pgfmathresult\commtikz@u\relax
    \fi
    \pgfkeysgetvalue{/tikz/commtikz/port span x}{\commtikz@psx}%
    \ifx\commtikz@psx\pgfutil@empty
      \commtikz@phw=\commtikz@hw\relax
    \else
      \pgfmathparse{\commtikz@psx}%
      \ifpgfmathunitsdeclared
        \commtikz@phw=0.5\dimexpr\pgfmathresult pt\relax
      \else
        \commtikz@phw=0.5\dimexpr\pgfmathresult\commtikz@u\relax
      \fi
    \fi
    \pgfkeysgetvalue{/tikz/commtikz/port span y}{\commtikz@psy}%
    \ifx\commtikz@psy\pgfutil@empty
      \commtikz@phh=\commtikz@hh\relax
    \else
      \pgfmathparse{\commtikz@psy}%
      \ifpgfmathunitsdeclared
        \commtikz@phh=0.5\dimexpr\pgfmathresult pt\relax
      \else
        \commtikz@phh=0.5\dimexpr\pgfmathresult\commtikz@u\relax
      \fi
    \fi
    \edef\commtikz@saved@u{\the\commtikz@u}%
    \commtikz@validatefixedports
    \commtikz@savearrowunit
    \edef\commtikz@saved@hw{\the\commtikz@hw}%
    \edef\commtikz@saved@hh{\the\commtikz@hh}%
    \edef\commtikz@saved@phw{\the\commtikz@phw}%
    \edef\commtikz@saved@phh{\the\commtikz@phh}%
    \edef\commtikz@numN{\pgfkeysvalueof{/tikz/commtikz/numN}}%
    \edef\commtikz@numS{\pgfkeysvalueof{/tikz/commtikz/numS}}%
    \edef\commtikz@numE{\pgfkeysvalueof{/tikz/commtikz/numE}}%
    \edef\commtikz@numW{\pgfkeysvalueof{/tikz/commtikz/numW}}%
    \edef\commtikz@portorder{\pgfkeysvalueof{/tikz/commtikz/port order}}%
    \addtosavedmacro\commtikz@saved@u
    \addtosavedmacro\commtikz@saved@hw
    \addtosavedmacro\commtikz@saved@hh
    \addtosavedmacro\commtikz@saved@phw
    \addtosavedmacro\commtikz@saved@phh
    \addtosavedmacro\commtikz@numN
    \addtosavedmacro\commtikz@numS
    \addtosavedmacro\commtikz@numE
    \addtosavedmacro\commtikz@numW
    \addtosavedmacro\commtikz@portorder
  }

  \savedanchor\centerpoint{\pgfpointorigin}
  \savedanchor\textpoint{%
    \pgf@x=-0.5\wd\pgfnodeparttextbox\relax
    \pgf@y=-0.5\ht\pgfnodeparttextbox\relax
    \advance\pgf@y by 0.5\dp\pgfnodeparttextbox\relax
  }

  \anchor{center}{\centerpoint}
  \anchor{text}{\textpoint}
  \anchor{mid}{\pgfpoint{0pt}{0pt}}
  \anchor{base}{\pgfpoint{0pt}{-0.5\ht\pgfnodeparttextbox}}

  \anchor{north}{\commtikz@rectparams\pgfpoint{0pt}{\commtikz@hh}}
  \anchor{south}{\commtikz@rectparams\pgfpoint{0pt}{-\commtikz@hh}}
  \anchor{east}{\commtikz@rectparams\pgfpoint{\commtikz@hw}{0pt}}
  \anchor{west}{\commtikz@rectparams\pgfpoint{-\commtikz@hw}{0pt}}
  \anchor{north east}{\commtikz@rectparams\pgfpoint{\commtikz@hw}{\commtikz@hh}}
  \anchor{north west}{\commtikz@rectparams\pgfpoint{-\commtikz@hw}{\commtikz@hh}}
  \anchor{south east}{\commtikz@rectparams\pgfpoint{\commtikz@hw}{-\commtikz@hh}}
  \anchor{south west}{\commtikz@rectparams\pgfpoint{-\commtikz@hw}{-\commtikz@hh}}

  % Short compass aliases
  \anchor{N}{\commtikz@rectparams\pgfpoint{0pt}{\commtikz@hh}}
  \anchor{S}{\commtikz@rectparams\pgfpoint{0pt}{-\commtikz@hh}}
  \anchor{E}{\commtikz@rectparams\pgfpoint{\commtikz@hw}{0pt}}
  \anchor{W}{\commtikz@rectparams\pgfpoint{-\commtikz@hw}{0pt}}
  \anchor{NE}{\commtikz@rectparams\pgfpoint{\commtikz@hw}{\commtikz@hh}}
  \anchor{NW}{\commtikz@rectparams\pgfpoint{-\commtikz@hw}{\commtikz@hh}}
  \anchor{SE}{\commtikz@rectparams\pgfpoint{\commtikz@hw}{-\commtikz@hh}}
  \anchor{SW}{\commtikz@rectparams\pgfpoint{-\commtikz@hw}{-\commtikz@hh}}
  \anchor{in}{\commtikz@rectparams\pgfpoint{-\commtikz@hw}{0pt}}
  \anchor{out}{\commtikz@rectparams\pgfpoint{\commtikz@hw}{0pt}}
  \anchor{input}{\commtikz@rectparams\pgfpoint{-\commtikz@hw}{0pt}}
  \anchor{output}{\commtikz@rectparams\pgfpoint{\commtikz@hw}{0pt}}

  % Dynamic multi-port anchors N0..N15, S0..S15, E0..E15, W0..W15
  \commtikz@defrectportanchors

  \anchorborder{%
    \pgf@xa=\pgf@x\relax
    \pgf@ya=\pgf@y\relax
    \commtikz@rectparams
    \pgfpointborderrectangle{\pgfpoint{\pgf@xa}{\pgf@ya}}{\pgfpoint{\commtikz@hw}{\commtikz@hh}}%
  }

  \backgroundpath{%
    \commtikz@rectparams
    \pgfpathrectanglecorners{\pgfpoint{-\commtikz@hw}{-\commtikz@hh}}{\pgfpoint{\commtikz@hw}{\commtikz@hh}}%
  }

  \beforebackgroundpath{%
    \commtikz@rectparams
    \commtikz@drawcomplabel{0pt}{0pt}%
  }
}

% ==============================================================================
% 4. Base Shape: commcircle
% ==============================================================================
\def\commtikz@circparams{%
  \commtikz@savedcirc
  \commtikz@arrowunit=\commtikz@saved@arrowunit\relax
  \commtikz@u=\commtikz@saved@u\relax
  \commtikz@r=\commtikz@saved@r\relax
  \commtikz@rk=\commtikz@saved@rk\relax
}

\pgfdeclareshape{commcircle}{
  \savedmacro\commtikz@savedlabel{\commtikz@savelabeldistance}
  \savedmacro\commtikz@savedcirc{%
    \commtikz@setunit
    \pgfmathparse{\pgfkeysvalueof{/tikz/commtikz/radius}}%
    \ifpgfmathunitsdeclared
      \commtikz@r=\pgfmathresult pt\relax
    \else
      \commtikz@r=\pgfmathresult\commtikz@u\relax
    \fi
    \commtikz@hw=\commtikz@r\relax
    \commtikz@hh=\commtikz@r\relax
    \commtikz@savearrowunit
    \commtikz@rk=0.70710678\commtikz@r\relax
    \edef\commtikz@saved@u{\the\commtikz@u}%
    \edef\commtikz@saved@r{\the\commtikz@r}%
    \edef\commtikz@saved@rk{\the\commtikz@rk}%
    \addtosavedmacro\commtikz@saved@u
    \addtosavedmacro\commtikz@saved@r
    \addtosavedmacro\commtikz@saved@rk
  }

  \savedanchor\centerpoint{\pgfpointorigin}
  \savedanchor\textpoint{%
    \pgf@x=-0.5\wd\pgfnodeparttextbox\relax
    \pgf@y=-0.5\ht\pgfnodeparttextbox\relax
    \advance\pgf@y by 0.5\dp\pgfnodeparttextbox\relax
  }

  \anchor{center}{\centerpoint}
  \anchor{text}{\textpoint}
  \anchor{mid}{\pgfpoint{0pt}{0pt}}
  \anchor{base}{\pgfpoint{0pt}{-0.5\ht\pgfnodeparttextbox}}

  \anchor{north}{\commtikz@circparams\pgfpoint{0pt}{\commtikz@r}}
  \anchor{south}{\commtikz@circparams\pgfpoint{0pt}{-\commtikz@r}}
  \anchor{east}{\commtikz@circparams\pgfpoint{\commtikz@r}{0pt}}
  \anchor{west}{\commtikz@circparams\pgfpoint{-\commtikz@r}{0pt}}
  \anchor{north east}{\commtikz@circparams\pgfpoint{\commtikz@rk}{\commtikz@rk}}
  \anchor{north west}{\commtikz@circparams\pgfpoint{-\commtikz@rk}{\commtikz@rk}}
  \anchor{south east}{\commtikz@circparams\pgfpoint{\commtikz@rk}{-\commtikz@rk}}
  \anchor{south west}{\commtikz@circparams\pgfpoint{-\commtikz@rk}{-\commtikz@rk}}

  % Aliases
  \anchor{N}{\commtikz@circparams\pgfpoint{0pt}{\commtikz@r}}
  \anchor{S}{\commtikz@circparams\pgfpoint{0pt}{-\commtikz@r}}
  \anchor{E}{\commtikz@circparams\pgfpoint{\commtikz@r}{0pt}}
  \anchor{W}{\commtikz@circparams\pgfpoint{-\commtikz@r}{0pt}}
  \anchor{N0}{\commtikz@circparams\pgfpoint{0pt}{\commtikz@r}}
  \anchor{S0}{\commtikz@circparams\pgfpoint{0pt}{-\commtikz@r}}
  \anchor{E0}{\commtikz@circparams\pgfpoint{\commtikz@r}{0pt}}
  \anchor{W0}{\commtikz@circparams\pgfpoint{-\commtikz@r}{0pt}}
  \anchor{NE}{\commtikz@circparams\pgfpoint{\commtikz@rk}{\commtikz@rk}}
  \anchor{NW}{\commtikz@circparams\pgfpoint{-\commtikz@rk}{\commtikz@rk}}
  \anchor{SE}{\commtikz@circparams\pgfpoint{\commtikz@rk}{-\commtikz@rk}}
  \anchor{SW}{\commtikz@circparams\pgfpoint{-\commtikz@rk}{-\commtikz@rk}}
  \anchor{in}{\commtikz@circparams\pgfpoint{-\commtikz@r}{0pt}}
  \anchor{out}{\commtikz@circparams\pgfpoint{\commtikz@r}{0pt}}
  \anchor{input}{\commtikz@circparams\pgfpoint{-\commtikz@r}{0pt}}
  \anchor{output}{\commtikz@circparams\pgfpoint{\commtikz@r}{0pt}}

  \anchorborder{%
    \pgf@xa=\pgf@x\relax
    \pgf@ya=\pgf@y\relax
    \commtikz@circparams
    % PGF's ellipse helper reuses xa/ya while reading the radii.
    \edef\commtikz@borderdirection{\noexpand\pgfpoint{\the\pgf@xa}{\the\pgf@ya}}%
    \pgfpointborderellipse{\commtikz@borderdirection}{\pgfpoint{\commtikz@r}{\commtikz@r}}%
  }

  \backgroundpath{%
    \commtikz@circparams
    \pgfpathcircle{\pgfpointorigin}{\commtikz@r}%
  }

  \beforebackgroundpath{%
    \commtikz@circparams
    \commtikz@drawcomplabel{0pt}{0pt}%
  }
}

\def\commtikz@inheritcircanchors#1{%
  \inheritsavedanchors[from=#1]%
  \inheritanchorborder[from=#1]%
  \inheritanchor[from=#1]{center}%
  \inheritanchor[from=#1]{text}%
  \inheritanchor[from=#1]{mid}%
  \inheritanchor[from=#1]{base}%
  \inheritanchor[from=#1]{north}%
  \inheritanchor[from=#1]{south}%
  \inheritanchor[from=#1]{east}%
  \inheritanchor[from=#1]{west}%
  \inheritanchor[from=#1]{north east}%
  \inheritanchor[from=#1]{north west}%
  \inheritanchor[from=#1]{south east}%
  \inheritanchor[from=#1]{south west}%
  \inheritanchor[from=#1]{N}%
  \inheritanchor[from=#1]{S}%
  \inheritanchor[from=#1]{E}%
  \inheritanchor[from=#1]{W}%
  \inheritanchor[from=#1]{N0}%
  \inheritanchor[from=#1]{S0}%
  \inheritanchor[from=#1]{E0}%
  \inheritanchor[from=#1]{W0}%
  \inheritanchor[from=#1]{NE}%
  \inheritanchor[from=#1]{NW}%
  \inheritanchor[from=#1]{SE}%
  \inheritanchor[from=#1]{SW}%
  \inheritanchor[from=#1]{in}%
  \inheritanchor[from=#1]{out}%
  \inheritanchor[from=#1]{input}%
  \inheritanchor[from=#1]{output}%
}

% Generic block styles
\tikzset{
  commrect/.style={shape=commrect, commtikz square node},
  rectblock/.style={commrect, #1},
  rectblock/.default={},
  commsquare/.style={shape=commrect, commtikz square node},
  squareblock/.style={commsquare, #1},
  squareblock/.default={},
  commcircle/.style={shape=commcircle, commtikz base node, comp radius=0.5},
  circleblock/.style={commcircle, #1},
  circleblock/.default={},
}
